%-------------------------------%
%  Author: Alessandro Sciarra   %
%    Date: 25 Jun 2019          %
%-------------------------------%

%~~~~~~~~~~~~~~~~~~~~~~~~~~~~~~~~~~~~~~~~~~~~%
\begin{frame}{Several kind of expansions}{\URL[PB]{https://www.gnu.org/software/bash/manual/}{Bash manual v5.0 section 3.5}}
    Expansion is performed on the command line after it has been split into tok\tikzmark{tokens}ens:\\[0.5em]
    \begin{itemize}
        \item Brace expansion\tikzmark{ExpA}
        \item Parameter and variable expansion\tikzmark{ExpB}
        \item Tilde expansion
        \item Arithmetic expansion
        \item Process substitution \tc{Gray!80}{\Remark{Not available if \URL[Gray]{https://www.gnu.org/software/bash/manual/html_node/Bash-POSIX-Mode.html}{Bash in POSIX Mode (item 30)}}}
        \item Command substitution\tikzmark{ExpC}
        \item Word splitting\tikzmark{ExpD}
        \item Filename expansion\tikzmark{ExpE}
    \end{itemize}
    \begin{varblock}{alerted}[0.85\textwidth]{Quote removal}
        After the preceding expansions, all unquoted occurrences of the characters\\
        \alert{\texttt{\textbackslash}}, \alert{\texttt{'}}, and \alert{\texttt{"}}
        that did not result from one of the above expansions are removed.
    \end{varblock}
    \begin{tikzpicture}[remember picture, overlay]
        \coordinate (xPos) at ($(current page.north)!0.3!(current page.north east)$);
        \draw[very thick, decorate, decoration={brace,amplitude=6pt}] (ExpB -| xPos) ++(-5mm,1mm) -- ($(ExpC -| xPos)+(-5mm,-1mm)$)
              coordinate[midway] (second);
        \node[anchor=west] at (second -| xPos) {\MakeEnumerateBox{2}$\quad$\alert{At the same time!}};
        \draw[from, shorter={0mm}{3mm}] (ExpA) -- (xPos |- ExpA) node[anchor=west] {\MakeEnumerateBox{1}};
        \draw[from, shorter={0mm}{3mm}] (ExpD) -- (xPos |- ExpD) node[anchor=west] {\MakeEnumerateBox{3}};
        \draw[from, shorter={0mm}{3mm}] (ExpE) -- (xPos |- ExpE) node[anchor=west] {\MakeEnumerateBox{4}};
        \node[anchor=north east, font=\scriptsize, above = of tokens] (metacharacters){
            \begin{tabular}{c}
                A line is split into tokens using unquoted metacharacters:\\[0.2em]
                \PP{\texttt{\textvisiblespace\ \textbackslash t \textbackslash n | \& ; ( ) < >}}
            \end{tabular}
        };
        \path[from, shorter={2mm}{0mm}] (tokens) edge[out=90, in=270] (metacharacters);
    \end{tikzpicture}
\end{frame}
%~~~~~~~~~~~~~~~~~~~~~~~~~~~~~~~~~~~~~~~~~~~~%
\begin{frame}[fragile]{Brace expansion}
    \vspace{-6mm}
    \begin{overlayarea}{\textwidth}{0.7\textheight}
        \begin{itemize}
            \item It comes in two forms:
                \begin{enumerate}
                    \item \bash|[prefix]{comma separated list}[postfix]|
                    \item \bash|[prefix]{X..Y[..increment]}[postfix]|
                \end{enumerate}
            \item<only@1> \bash|[prefix]| and \bash|[postfix]| are optional 
            \item<only@1> \bash|[prefix]| and \bash|[postfix]| might contain other brace expansions
            \item<only@1> Brace expansions can be nested and, if so, they work from the outside in
            \item<only@1> \bash|X| and \bash|Y| are either characters or numbers
            \item<only@1> \bash|increment| is an optional integer (if omitted, it is +1 or -1 as appropriate)
            \item<only@1> If \bash|X| and \bash|Y| are numbers, leading 0 are respected to force each term to have the same width
            \item<only@1> If the brace expansion syntax is not respected, then brace expansion is not performed!
        \end{itemize}
        \vspace{-3mm}
        \begin{varblock}{}[\textwidth]{Result of the expansion}<only@1>
            A space separated list of all combinations of \bash|[prefix]| and \bash|[postfix]| with the elements in the brace.
            In \MakeEnumerateBox{2}, the sequence in braces is at first completed going back to \MakeEnumerateBox{1}.
            Order is respected left to right. Multiple non-nested braces expand to all combinations.
        \end{varblock}
        \begin{onlyenv}<2>
            \begin{lstlisting}[style=MyBash, style=oddnumbers, aboveskip=5mm]
                $ echo {a,b,c}
                |+a b c+|
                $ echo {a,b,c}.tex
                |+a.tex b.tex c.tex+|
                $ echo image.{jpg,png,pdf}
                |+image.jpg image.png image.pdf+|
                $ echo {1..8}
                |+1 2 3 4 5 6 7 8+|
                $ echo {8..1}
                |+8 7 6 5 4 3 2 1+|
                $ echo {1..8..3}
                |+1 4 7+|
                $ echo {a..e}
                |+a b c d e+|
                $ echo {e..a..2}
                |+e c a+|
                $ echo file_{01..3}.pdf
                |+file_01.pdf file_02.pdf file_03.pdf+| #Note leading zeros!
            \end{lstlisting}
        \end{onlyenv}
        \begin{onlyenv}<3>
            \begin{lstlisting}[style=MyBash, style=oddnumbers, aboveskip=5mm, firstnumber=18]
                $ echo {A..C}{0..2}
                |+A0 A1 A2 B0 B1 B2 C0 C1 C2+|
                $ echo {A..C}@|\textvisiblespace|@{0..2} # Here two independent expansions!
                |+A B C 0 1 2+|
                $ echo {in,out}{go,com}ing
                |+ingoing incoming outgoing outcoming+|
                $ echo {{A,E,I,O,U},{0..9}}
                |+A E I O U 0 1 2 3 4 5 6 7 8 9+|
                $ echo {a..z..x}
                |+{a..z..x}+|  # No brace expansion!
                $ aVar=1; echo {$aVar..5}; unset aVar
                |+{1..5}+|     # No brace expansion!
                $ echo {b,1..5}
                |+b 1..5+|     # Not surprising, right?
                $ echo {b,{1..5}}
                |+b 1 2 3 4 5+|
                $ echo {b,{1..5}}} # Which of the last braces is kept?
                |+b} 1} 2} 3} 4} 5}+|
            \end{lstlisting}
        \end{onlyenv}
    \end{overlayarea}
\end{frame}
%~~~~~~~~~~~~~~~~~~~~~~~~~~~~~~~~~~~~~~~~~~~~%
\begin{frame}{Parameter expansion: Overview}
    \vspace{-3mm}
    \begin{itemize}
        \item It is the most used form of expansion!
        \item Its basic syntax consists of a single dollar sign alone
        \item Braces after the dollar are only sometimes necessary\\
              $\to$ it is good practice to always use them: \PP{\texttt{\$\{\ldots\}}}
        \item Basic form: \PB{\texttt{\$\{parameter\}}} $\,\to\,$ The value of parameter is substituted
        \item There are plenty of incredibly powerful modified forms of parameter expansion\\
              $\to$ you will learn them by using, but \alert{keep in mind they exist}!
        \item General form: \PB{\texttt{\$\{[prefix]parameter[postfix]\}}}
        \item Not all (particular) cases are discussed here immediately
        \item More on parameter expansion in the future\\
              \begin{itemize}
                  \item[$\to$] namerefs and indirect expansion
                  \item[$\to$] arrays and associative arrays
              \end{itemize}
    \end{itemize}
    \FrameRemark{Everything can be found in the \URL[PB]{https://www.gnu.org/software/bash/manual/}{Bash manual v5.0 section 3.5.3}}
\end{frame}
%~~~~~~~~~~~~~~~~~~~~~~~~~~~~~~~~~~~~~~~~~~~~%
\begin{frame}[fragile]{Parameter expansion: Checking variables state}
    \vspace{-3mm}
    \begin{onlyenv}<1>
        \begin{varblock}{alerted}[0.8\textwidth]{}
            \small\alert{In each of the cases below, word is subject to tilde expansion, parameter expansion, command substitution, and arithmetic expansion!}
        \end{varblock}
        \begin{center}
            \begin{tabular}{r@{\quad}>{\footnotesize}l}
                \makecell[tc]{
                    \PB{\small\texttt{\$\{parameter:-word\}}}\\[-0.5em]
                    \PP{\ssmall\textbf{Use Default Value}}
                } & \makecell[tl]{Parameter unset or null: \PP{expansion of word is substituted}\\
                                  Otherwise: the value of parameter is substituted} \\
                \makecell[tc]{
                    \PB{\small\texttt{\$\{parameter:=word\}}}\\[-0.5em]
                    \PP{\ssmall\textbf{Assign Default Value}}
                } & \makecell[tl]{Parameter unset or null: \PP{expansion of word is assigned to parameter}\\
                                  Otherwise: the value of parameter remains unchanged\\
                                  The value of parameter is then substituted\\
                                  \Remark[0pt]{Positional parameters and special parameters may not be assigned to in this way}}\\
                \makecell[tc]{
                    \PB{\small\texttt{\$\{parameter:?word\}}}\\[-0.5em]
                    \PP{\ssmall\textbf{Exit with message}}
                } & \makecell[tl]{Parameter is null or unset: \PP{the expansion of word to the standard error}\\
                                  The shell, then, if it is not interactive, exits\\
                                  Otherwise, the value of parameter is substituted\\
                                  \Remark[0pt]{If word is not present, a default message is produced}}\\
                \makecell[tc]{
                    \PB{\small\texttt{\$\{parameter:+word\}}}\\[-0.5em]
                    \PP{\ssmall\textbf{Use Alternative value}}
                } & \makecell[tl]{Parameter is null or unset: \PP{nothing is substituted}\\
                                  Otherwise, the expansion of word is substituted}\\
            \end{tabular}
        \end{center}
        \vspace{-3mm}
        \begin{varblock}{}[0.8\textwidth]{A different variant}
            \small Omitting the colon results in a test only for a parameter that is unset. \\
            Put another way, if the colon is omitted, the operator tests only for existence.
        \end{varblock}
    \end{onlyenv}
    \begin{onlyenv}<2>
        \begin{lstlisting}[style=MyBash, style=oddnumbers, style=smaller]
            $ aVar="Hello"
            # The variable aVar is set and not empty/null
            $ echo "_${aVar}_"
            |+_Hello_+|
            $ echo "_${aVar-Goodbye}_   _${aVar:-Goodbye}_"  # Use default Value
            |+_Hello_   _Hello_+|
            $ echo "_${aVar=Goodbye}_   _${aVar:=Goodbye}_"  # Assign default Value
            |+_Hello_   _Hello_+|
            $ echo "_${aVar?Goodbye}_   _${aVar:?Goodbye}_"  # Exit with message
            |+_Hello_   _Hello_+|
            $ echo "_${aVar+Goodbye}_   _${aVar:+Goodbye}_"  # Use alternative value
            |+_Goodbye_   _Goodbye_+|
            $ aVar=""
            # The variable aVar is now set BUT empty/null
            $ echo "_${aVar}_"
            |+__+|
            $ echo "_${aVar-Goodbye}_   _${aVar:-Goodbye}_"
            |+__   _Goodbye_+|
            $ echo "_${aVar=Goodbye}_   _${aVar:=Goodbye}_"; aVar=""
            |+__   _Goodbye_+|
            $ echo "_${aVar?Goodbye}_   _${aVar:?Goodbye}_"
            |+bash: aVar: Goodbye+| # Shell STOPS at second expansion and prints message!
            $ echo "_${aVar?Goodbye}_"
            |+__+|
            $ echo "_${aVar+Goodbye}_   _${aVar:+Goodbye}_"
            |+_Goodbye_   __+|
        \end{lstlisting}
    \end{onlyenv}
    \begin{onlyenv}<3>
        \begin{lstlisting}[style=MyBash, style=oddnumbers, style=smaller, firstnumber=26]
            $ unset aVar
            # The variable aVar is now UNSET
            $ echo "_${aVar}_"
            |+__+|
            $ echo "_${aVar-Goodbye}_   _${aVar:-Goodbye}_"
            |+_Goodbye_   _Goodbye_+|
            $ echo "_${aVar=Goodbye}_"; unset aVar
            |+_Goodbye_+|
            $ echo "_${aVar:=Goodbye}_"; unset aVar
            |+_Goodbye_+|
            $ echo "_${aVar?Goodbye}_"
            |+bash: aVar: Goodbye+| # Shell STOPS at expansion and prints message!
            $ echo "_${aVar:?Goodbye}_"
            |+bash: aVar: Goodbye+| # Shell STOPS at expansion and prints message!
            $ echo "_${aVar+Goodbye}_   _${aVar:+Goodbye}_"
            |+__   __+|
        \end{lstlisting}
        \begin{varblock}{}[0.9\textwidth]{Ok, but is all this useful? When?}
            Ohhh yes! These kinds of expansion are useful e.g. when it comes to checking whether a variable is set, unset or set but empty.
            Sometimes there are alternatives, but sometimes you need exactly this!
        \end{varblock}
    \end{onlyenv}
\end{frame}
%~~~~~~~~~~~~~~~~~~~~~~~~~~~~~~~~~~~~~~~~~~~~%
\begin{frame}[fragile]{Parameter expansion: String manipulation (I)}
    \vspace{-3mm}
    \setbeamerfont{description item}{series=\bfseries}
    \setbeamercolor{description item}{fg=PP}
    \begin{description}[Substring Expansion:]
        \item[String Length:] \PB{\small\texttt{\$\{\#parameter\}}}\\
            {\small
                The length in characters of the value of parameter is substituted
            }
        \item[Substring Expansion:] \PB{\small\texttt{\$\{parameter:offset:length\}}}\\
            {\small
                Expands to up to \texttt{length} characters of \texttt{parameter} starting at the character specified by \texttt{offset} (0-indexed)\\
                If \texttt{:length} is omitted, go all the way to the end\\[-0.5em]
                \Remark[0pt]{If \texttt{offset} is negative, count backward from the end of \texttt{parameter} instead of forward from the beginning}
            }
    \end{description}
    \begin{uncoverenv}<2>
        \begin{lstlisting}[style=MyBash, style=oddnumbers, aboveskip=2mm]
            $ aVar='01234567890abcdefgh'
            # Let us demonstrate a bit
            $ echo "String ${aVar} is ${#aVar} characters long"
            |+String 01234567890abcdefgh is 19 characters long+|
            $ echo "${aVar:7} and 0 length: \"${aVar:7:0}\""
            |+7890abcdefgh and 0 length: ""+|
            $ echo "${aVar:7:2} and ${aVar:7:-2}"
            |+78 and 7890abcdef+|
            $ echo "${aVar:@|\textvisiblespace|@-7}" # Why do you need a space?
            |+bcdefgh+|
            $ echo "${aVar: -7:2} and ${aVar: -7:0} and ${aVar: -7:-2}"
            |+bc and  and bcdef+|
        \end{lstlisting}
    \end{uncoverenv}
    \FrameRemark{These types of expansion apply also to arrays: We will come back to this!}
\end{frame}
%~~~~~~~~~~~~~~~~~~~~~~~~~~~~~~~~~~~~~~~~~~~~%
\begin{frame}[fragile]{Parameter expansion: String manipulation (II)}
    \vspace{-3mm}
    \setbeamerfont{description item}{series=\bfseries}
    \setbeamercolor{description item}{fg=PP}
    \begin{description}[Remove Smallest Suffix:]
        \item[Remove Smallest Prefix:] \PB{\small\texttt{\$\{parameter\#pattern\}}}\\
            {\small
                The \texttt{pattern} is matched against the \textbf{beginning} of \texttt{parameter}\\
                which is expanded with the \textbf{shortest} match deleted
            }
        \item[Remove Largest Prefix:] \PB{\small\texttt{\$\{parameter\#\#pattern\}}}
            {\scriptsize
                As the previous, but the \textbf{longest} match is deleted
            }
        \item[Remove Smallest Suffix:] \PB{\small\texttt{\$\{parameter\%pattern\}}}\\
            {\small
                The \PB{\texttt{pattern}} is matched against the \textbf{end} of \texttt{parameter}\\
                which is expanded with the \textbf{shortest} match deleted
            }
        \item[Remove Largest Suffix:] \PB{\small\texttt{\$\{parameter\%\%pattern\}}}
            {\scriptsize
                As the previous, but the \textbf{longest} match is deleted
            }
    \end{description}
    \begin{uncoverenv}<2>
        \begin{lstlisting}[style=MyBash, style=oddnumbers, aboveskip=2mm]
            $ aVar='b1.2345_s0000_s1111'
            # Let us demonstrate a bit (* matches any characters)
            $ echo "${aVar#*_} and ${aVar##*_}"
            |+s0000_s1111 and s1111+|
            $ echo "${aVar%_*} and ${aVar%%_*}"
            |+b1.2345_s0000 and b1.2345+|
            $ echo "${aVar#NoMatch} and ${aVar%NoMatch}"
            |+b1.2345_s0000_s1111 and b1.2345_s0000_s1111+|
        \end{lstlisting}
    \end{uncoverenv}
    \FrameRemark{These types of expansion apply also to arrays: We will come back to this!}
\end{frame}
%~~~~~~~~~~~~~~~~~~~~~~~~~~~~~~~~~~~~~~~~~~~~%
\begin{frame}[fragile]{Parameter expansion: String manipulation (III)}
    \vspace{-3mm}
    \setbeamerfont{description item}{series=\bfseries}
    \setbeamercolor{description item}{fg=PP}
    \begin{description}[Replace at start:]
        \item[Replace first:] \PB{\small\texttt{\$\{parameter/pattern/string\}}}\\
            {\small
                Results in the expanded value of \texttt{parameter} with the first (unanchored) match\\
                of \texttt{pattern} replaced by \texttt{string}\\[-0.5em]
                \Remark[0pt]{Assume null string when the \texttt{/string} part is absent}
            }
        \item[Replace all:] \PB{\small\texttt{\$\{parameter//pattern/string\}}}\\
            {\small
                As the previous, but \textbf{every} match is replaced
            }
        \item[Replace at start:] \PB{\small\texttt{\$\{parameter/\#pattern/string\}}}\\
            {\small
                As the first, but matched against the \textbf{beginning}
            }
        \item[Replace at end:] \PB{\small\texttt{\$\{parameter/\%pattern/string\}}}\\
            {\small
                As the first, but matched against the \textbf{end}
            }
    \end{description}
    \begin{uncoverenv}<2>
        \begin{lstlisting}[style=MyBash, style=oddnumbers, aboveskip=2mm]
            $ aVar='123aa000aa321'
            # Let us demonstrate a bit
            $ echo "${aVar/aa/_} and ${aVar//aa/_}"
            |+123_000aa321 and 123_000_321+|
            $ echo "${aVar/#1/_} and ${aVar/%1/_}"
            |+_23aa000aa321 and 123aa000aa32_+|
            $ echo "${aVar/#/prefix} and ${aVar/%/postfix}"
            |+prefix123aa000aa321 and 123aa000aa321postfix+|
        \end{lstlisting}
    \end{uncoverenv}
    \FrameRemark{These types of expansion apply also to arrays: We will come back to this!}
\end{frame}
%~~~~~~~~~~~~~~~~~~~~~~~~~~~~~~~~~~~~~~~~~~~~%
\begin{frame}[fragile]{Parameter expansion: String manipulation (IV)}{Since Bash v4.0 (2009)}
    \vspace{-3mm}
    \setbeamerfont{description item}{series=\bfseries}
    \setbeamercolor{description item}{fg=PP}
    \begin{description}[Replace at start:]
        \item[First uppercase:] \PB{\small\texttt{\$\{parameter\^{}\}}}\\
            {\small
                Results in the expanded value of \texttt{parameter} with the first character capitalised
            }
        \item[All uppercase:] \PB{\small\texttt{\$\{parameter\^{}\^{}\}}}\\
            {\small
                As the previous, but all characters are capitalised
            }
        \item[First uppercase:] \PB{\small\texttt{\$\{parameter,\}}}\\
            {\small
                Results in the expanded value of \texttt{parameter} with the first character to lowercase
            }
        \item[All lowercase:] \PB{\small\texttt{\$\{parameter,,\}}}\\
            {\small
                As the previous, but all characters are transformed to lowercase
            }
    \end{description}
    \begin{uncoverenv}<2>
        \begin{lstlisting}[style=MyBash, style=oddnumbers, aboveskip=2mm]
            $ aVar='hELLO, World!'
            # Let us demonstrate a bit
            $ echo "${aVar^}  -  ${aVar^^}"
            |+HELLO, World!  -  HELLO, WORLD!+|
            $ echo "${aVar,}  -  ${aVar,,}"
            |+hELLO, World!  -  hello, world!+|
        \end{lstlisting}
    \end{uncoverenv}
    \FrameRemark{These types of expansion apply also to arrays: We will come back to this!}
\end{frame}
%~~~~~~~~~~~~~~~~~~~~~~~~~~~~~~~~~~~~~~~~~~~~%
\begin{frame}{Parameter expansion: String manipulation summary}
    \vspace{-3mm}
    \begin{varblock}{alerted}[0.88\textwidth]{\textbf{Good practice:}}
        You may be tempted to use external applications such as \bash|sed|, \bash|awk|, \bash|cut|, \bash|perl| or others to modify your strings.
        Be aware that all of these require an extra process to be started, which in some cases can cause slowdowns.\\
        \alert{Parameter Expansions are the perfect alternative!}
    \end{varblock}
    \bigskip
    \centerline{\includegraphics[width=0.5\textwidth, clip, trim=0 5mm 0 0]{GotPower}}
\end{frame}
%~~~~~~~~~~~~~~~~~~~~~~~~~~~~~~~~~~~~~~~~~~~~%
\begin{frame}[fragile]{Tilde expansion in its basic version}
    \vspace{-2mm}
    \begin{enumerate}
        \small
        \item \PP{If a word begins with an unquoted tilde character \texttt{\textasciitilde}}\\
              $\;\Rightarrow\;$ all of the characters up to the first unquoted slash are considered a \textbf{\PB{tilde-prefix}}.\\
              \phantom{$\;\Rightarrow\;$} \Remark[0pt]{all characters are considered, if there is no unquoted slash}
        \item \PP{If none of the characters in the tilde-prefix are quoted}\\
              $\;\Rightarrow\;$ the characters in the tilde-prefix following the tilde are treated as a possible \textbf{\PS{login name}}.
        \item \PP{If this login name is the null string}\\
              $\;\Rightarrow\;$ the tilde is replaced with the value of the \bash|HOME| shell variable,\\
              \phantom{$\;\Rightarrow\;$} or with the home directory of the user executing the shell, if \bash|HOME| is unset.\\
              \PP{Otherwise}\\
              $\;\Rightarrow\;$ the \textbf{\PB{tilde-prefix}} is replaced with the home directory associated with the specified \textbf{\PS{login name}}.
        \item \alert{If the login name is invalid the word is left unchanged.}
    \end{enumerate}
    \begin{uncoverenv}<2>
        \begin{lstlisting}[style=MyBash, style=oddnumbers, aboveskip=4mm]
            $ echo ~  ~palao
            |+/home/sciarra /home/palao+|
            $ echo ~"sciarra" ~sciar"ra" ~sciarra""
            |+~sciarra ~sciarra ~sciarra+|
            $ echo ~NotExistentUser # Invalid login name
            |+~NotExistentUser+|
        \end{lstlisting}
    \end{uncoverenv}
    \FrameRemark{Each variable assignment is also checked for unquoted tilde-prefixes immediately following a \bash|':'| or the first \bash|'='|.  In these cases, tilde expansion is also performed.}
\end{frame}
%~~~~~~~~~~~~~~~~~~~~~~~~~~~~~~~~~~~~~~~~~~~~%
\begin{frame}{The Directory Stack}{\URL[PB]{https://www.gnu.org/software/bash/manual/}{Bash manual v5.0 section 6.8}}
    \vspace{-3mm}
    \begin{columns}[c]
        \begin{column}{0.1\textwidth}
        \end{column}
        \begin{column}{0.6\textwidth}
            \begin{varblock}{example}[0.9\textwidth]{Intuitive definition of stack}
                An abstract data type or data structure based on the principle of last in first out.
            \end{varblock}
        \end{column}
        \begin{column}{0.3\textwidth}
        \end{column}
    \end{columns}
    \vspace{6mm}
    \begin{itemize}
        \item The directory stack is a list of recently-visited directories
        \item The \bash|dirs| builtin \PP{\textbf{displays}} the contents of the directory stack \\[1mm]
        \item \alert{The current directory is always the ``top'' of the directory stack}
        \item The \bash|pushd| builtin \PS{\textbf{adds}} directories to the  stack as it changes the current directory
        \item The \bash|popd| builtin \PT{\textbf{removes}} specified directories from the stack and changes the current directory to the directory removed
        \item Both \bash|pushd| and \bash|popd| accept a \PB{\texttt{-n}} option to avoid changing directory
    \end{itemize}
    \begin{tikzpicture}[remember picture, overlay]
        \foreach \l [count=\n from 0, evaluate=\l as \shade using \l*20] in {5,4,...,0}{
            \coordinate (tmpCoord) at ($(current page.east)+(-25mm,\n*6mm)$);
            \colorlet{tmpColor}{PP!\shade!PS}
            \node[draw, minimum size=6mm, text=tmpColor, fill=tmpColor!10] at (tmpCoord) {\l};
        }
        \path[from] ($(tmpCoord)+(-1mm,3mm)$) edge[out=90, in=0] node[pos=1, left] {Last In}   ++(-5mm,5mm);
        \path[to] ($(tmpCoord)+(+1mm,3mm)$) edge[out=90, in=180] node[pos=1, right] {First Out} ++(+5mm,5mm);
    \end{tikzpicture}
    \FrameRemark{The contents of the directory stack are also visible as the value of the \,\bash|DIRSTACK|\, shell array variable.}
\end{frame}
%~~~~~~~~~~~~~~~~~~~~~~~~~~~~~~~~~~~~~~~~~~~~%
\begin{frame}[fragile]{The Directory Stack: An example usage}
    \vspace{-3mm}
    \begin{lstlisting}[style=MyBash]
        $ cd ~; pwd; dirs -v
        |+/home/sciarra+|
        |+ 0  ~+|
        $ cd ~/Documents/Seminars/BashCourse_2020; dirs -v
        |+ 0  ~/Documents/Seminars/BashCourse_2020+|
        $ pushd /myBackupDisk/University/Seminars/BashCourse_2020
        # Stack automatically displayed on one line
        $ dirs -v
        |+ 0  /myBackupDisk/University/Seminars/BashCourse_2020+|
        |+ 1  ~/Documents/Seminars/BashCourse_2020+|
        $ pushd ~; pwd
        # Stack automatically displayed on one line
        |+/home/sciarra+|
        $ dirs -v
        |+ 0  ~+|
        |+ 1  /myBackupDisk/University/Seminars/BashCourse_2020+|
        |+ 2  ~/Documents/Seminars/BashCourse_2020+|
        $ cd Documents; dirs -v
        |+ 0  ~/Documents+|
        |+ 1  /myBackupDisk/University/Seminars/BashCourse_2020+|
        |+ 2  ~/Documents/Seminars/BashCourse_2020+|
        $ popd; pwd  # Stack automatically displayed on one line
        |+/myBackupDisk/University/Seminars/BashCourse_2020+|
    \end{lstlisting}
\end{frame}
%~~~~~~~~~~~~~~~~~~~~~~~~~~~~~~~~~~~~~~~~~~~~%
\begin{frame}[fragile]{Tilde expansion and its advanced usage}
    \vspace*{-3mm}
    \begin{overlayarea}{\textwidth}{0.65\textheight}
        \begin{onlyenv}<1-2>
            \PB{If the tilde-prefix matches the following forms, a special expansion is done:}
            \vspace{-4mm}
            \begin{columns}
                \begin{column}{\dimexpr\paperwidth-10mm}
                    \begin{description}%[\texttt{1 2 \ldots}]
                        \setlength{\itemsep}{2mm}
                        \item[\texttt{\textasciitilde+}]
                            \texttt{\$\{\tc{environment-color}{PWD}\}} is substituted.
                        \item[\texttt{\textasciitilde-}]
                            The value of the shell variable \bash|OLDPWD|, if it is set, is substituted.\\
                            $\;\to\;$ \texttt{\$\{\tc{environment-color}{OLDPWD}-\tc{strings-color}{'\textasciitilde-'}\}}
                        \item[\texttt{\textasciitilde N}]
                            The string that would be displayed by \enspace\bash|dirs +N|\enspace is substituted.
                        \item[\texttt{\textasciitilde+N}]
                            The string that would be displayed by \enspace\bash|dirs +N|\enspace is substituted.
                        \item[\texttt{\textasciitilde-N}]
                            The string that would be displayed by \enspace\bash|dirs -N|\enspace is substituted.
                    \end{description}
                \end{column}
            \end{columns}
            \vspace{2mm}
            \begin{varblock}{alert}[0.68\textwidth]{}
                \large\alert{If the tilde expansion fails, the word is left unchanged!}
            \end{varblock}
        \end{onlyenv}
        \begin{onlyenv}<3->
            \begin{lstlisting}[style=MyBash]
                $ dirs -v
                |+ 0  ~+|
                |+ 1  /scratch/latticeqcd/sciarra+|
                |+ 2  /home/latticeqcd/public+|
                |+ 3  /arc01/archive/latticeqcd/sciarra+|
                |+ 4  /arc02/archive/latticeqcd/sciarra+|
                $ cd ~2; pwd
                |+/home/latticeqcd/public+|
                $ cd ~-0; pwd; cd ~4; pwd
                |+/arc02/archive/latticeqcd/sciarra+|
                |+/arc02/archive/latticeqcd/sciarra+|
                # Of course you can use them for any operation, e.g.
                $ cp -r ~3/Results ~4/Results_backup
                $ ls ~5
                |+ls: cannot access ~5: No such file or directory+|
            \end{lstlisting}
        \end{onlyenv}
    \end{overlayarea}
    \begin{varblock}{example}[0.65\textwidth]{}<uncover@2->
        \large\PS{You can then take advantage of the Directory Stack!}
    \end{varblock}
    \FrameRemark{Bash also performs tilde expansion on words satisfying the conditions of variable assignments (e.g. \bash|alias|, \bash|declare|, \bash|export|, \bash|readonly|, \bash|local|).}
\end{frame}
%~~~~~~~~~~~~~~~~~~~~~~~~~~~~~~~~~~~~~~~~~~~~%
\begin{frame}[fragile]{Arithmetic expansion}
    \vspace{-3mm}
    \begin{itemize}[<only@1>]
        \item Arithmetic expansion syntax is \PP{\texttt{\$((\ldots))}} and it can be nested
        \item \alert{Bash is only capable of integer arithmetic}
        \item If you need to do \textbf{a lot} of non-integer arithmetic, then Bash is the wrong tool!
        \item Shell variables are allowed as operands\\ $\to\,$ parameter expansion is performed before the expression is evaluated!
        \item Within an expression, shell variables may also be referenced by name without \PP{\texttt{\$\{\ldots\}}}\\
              $\Rightarrow\,$ strings are recursively interpreted as variable name until an integer value is found!
        \item A shell variable that is null or unset evaluates to
              \begin{itemize}
                  \item[$\circ$] \alert{\texttt{0}} when referenced by name \alert{\textbf{without using}} the parameter expansion syntax
                  \item[$\circ$] \PP{\texttt{''}} when referenced by name \PP{\textbf{using}} the parameter expansion syntax
              \end{itemize}
        \item Overflow is not checked but division by 0 gives an error
        \item Operators and their precedence are the same as in the \texttt{C} language\\ $\to\,$ \URL[PB]{https://www.gnu.org/software/bash/manual/}{Bash manual v5.0 section 6.5}
    \end{itemize}
    \begin{onlyenv}<2>
        \begin{lstlisting}[style=MyBash, style=oddnumbers, style=smaller]
            $ echo $(( 22 / 3 ))
            |+7+|
            $ aVar=22; echo $(( ${aVar} / 3 )); unset aVar
            |+7+|
            $ aVar=22; echo $(( aVar / 3 )); unset aVar
            |+7+|
            $ echo $(( aVar / 3 ))     # aVar is unset,  no  ${} => 0 is used:  $((0/3))
            |+0+|
            $ echo $(( ${aVar} / 3 ))  # aVar is unset, with ${} => '' is used: $(( /3))
            |+bash: / 3 : syntax error: operand expected (error token is "/ 3 ")+|
            $ echo $(( ))
            |+0+|
            $ aVar='Hello'; echo $(( aVar )) $(( ${aVar} )); unset aVar
            |+0 0+| # No variable 'Hello' existing
            $ aVar='bVar'; bVar=5; echo $(( aVar )) $(( ${aVar} )); unset aVar bVar
            |+5 5+| # aVar is expanded to bVar, which is defined to 5
            $ aVar='HelloMyLove'; echo $(( aVar )) $(( ${aVar} )); unset aVar
            |+0 0+| # No variable 'HelloMyLove' existing
            $ aVar='Hello <3'; echo $(( aVar )) $(( ${aVar} )); unset aVar
            |+1 1+| # It makes sense, doesn't it?
            $ echo $(( 2**51 * 4096 ))
            |+-9223372036854775808+| # No error, no overflow check!
        \end{lstlisting}
    \end{onlyenv}
\end{frame}
%~~~~~~~~~~~~~~~~~~~~~~~~~~~~~~~~~~~~~~~~~~~~%
\begin{frame}[fragile]{Word splitting}{\ldots{}the evil of Bash!}
    \vspace{-6mm}
    \begin{onlyenv}<1>
        \begin{varblock}{}[0.9\textwidth]{When does it occur?}
            The shell scans the results of parameter expansion, command substitution, and arithmetic expansion that did not occur within double quotes for word splitting!
        \end{varblock}
        \begin{itemize}
            \item The shell treats each character of \bash|IFS| as a delimiter, and splits the results of the other expansions into words using these characters \textbf{as field \alert{\underline{terminators}}} (these are all dropped)
            \item Whitespace characters $\;$\PP{\texttt{\textvisiblespace\textbackslash{}t\textbackslash{}n}}$\;$ contained in \bash|IFS| are called \PP{IFS whitespaces}
            \item \alert{IFS whitespaces at the beginning/end} of the results of previous expansions \alert{are ignored}
            \item \alert{Adjacent IFS whitespaces} are considered all together as \alert{single field terminator}
            \item \alert{Any other character} in \bash|IFS| that is not an IFS whitespace \alert{\textbf{always} delimits a field}
            \item If the value of \bash|IFS| is null, no word splitting occurs
            \item Implicit null arguments, resulting from the expansion of empty parameters, are removed
            \item \PQ{\textbf{Note that if no expansion occurs, no splitting is performed}}
        \end{itemize}
    \end{onlyenv}
    \begin{onlyenv}<2>
        \begin{lstlisting}[style=MyBash, style=oddnumbers, aboveskip=3mm]
            # The args command puts the arguments it gets in <...>
            # NOT standard, I implemented it!
            $ args Hello world "I am Alessandro"
            |+3 args: <Hello> <world> <I am Alessandro>+|
            $ aVar="This is a variable"; args ${aVar}
            |+4 args: <This> <is> <a> <variable>+|
            $ args "${aVar}"
            1 args: <This is a variable>
            $ ls
            Day1.tex   Day2.tex   Day3.tex
            $ args $(ls)
            |+3 args: <Day1.tex> <Day2.tex> <Day3.tex>+|
            $ aVar="     This is a variable   "; args ${aVar}
            |+4 args: <This> <is> <a> <variable>+|
            $ aVar="     This   is  a    variable   "; args ${aVar}
            |+4 args: <This> <is> <a> <variable>+|
            $ aVar="/home/sciarra/Documents/Seminars/"; args ${aVar}
            |+1 args: </home/sciarra/Documents/Seminars/>+|
            $ IFS='/'; args ${aVar}; unset IFS
            |+5 args: <> <home> <sciarra> <Documents> <Seminars>+|
            # Why is there an empty argument at the beginning?
            # Why isn't there one at the end?
        \end{lstlisting}
    \end{onlyenv}
    \begin{onlyenv}<3>
        \begin{lstlisting}[style=MyBash, style=oddnumbers, aboveskip=3mm, firstnumber=22]
            # Let us have a look to more complex situations
            $ aVar="   /home/sciarra/Documents/Seminars   "
            $ IFS='/'; args ${aVar}; unset IFS
            |+5 args: <   > <home> <sciarra> <Documents> <Seminars   >+|
            $ IFS='@|\textvisiblespace|@/'; args ${aVar}; unset IFS # Note the space in IFS
            |+5 args: <> <home> <sciarra> <Documents> <Seminars>+|
            # Remember $'' special quoting to allow backslash escaping
            $ aVar=$'\n   \n  xxx \t   \t  xx  \t'
            $ args ${aVar}
            |+2 args: <xxx> <xx>+|
            $ IFS=''; args ${aVar}; unset IFS
            |+1 args: <+|
            |+   +|
            |+  xxx         xx    >+|
            $ IFS=$'\n'; args ${aVar}; unset IFS
            |+2 args: <   > <  xxx        xx    >+|
            $ IFS=$'\n\t'; args ${aVar}; unset IFS
            |+4 args: <   > <  xxx > <   > <  xx  >+|
            $ IFS=$'\n\t x'; args ${aVar}; unset IFS
            |+5 args: <> <> <> <> <>+|
            $ IFS=':'; args aaa:bbb; unset IFS aVar
            |+1 args: <aaa:bbb>+| # It makes sense, right?
        \end{lstlisting}
    \end{onlyenv}
\end{frame}
%~~~~~~~~~~~~~~~~~~~~~~~~~~~~~~~~~~~~~~~~~~~~%
\begin{frame}[fragile]{Word splitting: conclusion}
    \vspace{-4  mm}
        \begin{varblock*}{}[0.9\textwidth]{When does NOT it occur?}
        Word splitting is not performed
        \begin{itemize}
            \item on expansions inside Bash keywords such as \bash|[[| \texttt{\ldots} \bash|]]| and \bash|case|;
            \item on expansions in assignments. \\[0.3em] Thus, one does not need to quote anything in a command like these:
                  \begin{lstlisting}[style=MyBash, numbers=none, xrightmargin=45mm, belowskip=-5mm, aboveskip=2mm]
                      aVar=${bVar}; aVar=$(date)
                      PATH=/usr/local/bin:${PATH} 
                  \end{lstlisting}
                  Quoting anyway will not break anything, so \PP{if in doubt, quote!}
        \end{itemize}
    \end{varblock*}
    \begin{varblock}{alerted}[0.98\textwidth]{Remember:}
        \alert{Word Splitting is performed \textbf{after} Parameter Expansion has been performed!} \\
        That means that it is absolutely vital that we quote our Parameter Expansions, in case they may expand values that contain syntactical whitespace which will then be word-split.
    \end{varblock}
    \FrameRemark{It is almost never desirable to put syntactical whitespace in parameters.}
\end{frame}
%~~~~~~~~~~~~~~~~~~~~~~~~~~~~~~~~~~~~~~~~~~~~%
\begin{frame}{Other expansions}
    \vspace{-8mm}
    \begin{columns}
        \begin{column}{\dimexpr\paperwidth-10mm}
            \begin{center}
                \large\fbox{\alert{\textbf{Before}} Word Splitting:}
            \end{center}
            \begin{description}[Command substitution:]
                \setlength{\itemsep}{3mm}
                \item[Command substitution:]
                    It occurs when a command is enclosed as \alert{\texttt{\$(\ldots)}}\\
                    Do not use the deprecated syntax \tc{Gray!30}{\texttt{\textasciigrave\ldots\textasciigrave}}\\
                    Still to early to be discussed $\to$ \URL[PB]{https://www.gnu.org/software/bash/manual/}{Bash manual v5.0 section 3.5.4}\\
                    \Remark[0pt]{\alert{\texttt{\$(\ldots)}} is substituted by the standard output of the command (executed in a subshell), with any trailing newlines deleted.}
                \item[Process substitution:]
                    Useful, but not discussed here $\to$ \URL[PB]{https://www.gnu.org/software/bash/manual/}{Bash manual v5.0 section 3.5.6}\\
                    It takes the form \alert{\texttt{<(\ldots)}} or \alert{\texttt{>(\ldots)}}\\
                    It will discussed in the section about redirection
            \end{description}
            \vspace{2mm}
            \begin{center}
                \large\fbox{\PP{\textbf{After}} Word Splitting:}
            \end{center}
            \begin{description}[Command substitution:]%[\hspace{0.4\textwidth}]
                \item[Filename expansion:]
                    Discussed in the next section about globbing
            \end{description}
        \end{column}
    \end{columns}
\end{frame}
%~~~~~~~~~~~~~~~~~~~~~~~~~~~~~~~~~~~~~~~~~~~~%
